% Generated draft: requires mathematical and source review.
\documentclass[11pt]{article}
\usepackage[margin=25mm]{geometry}
\usepackage{fontspec}
\setmainfont{Latin Modern Roman}
\usepackage{amsmath,amssymb,mathtools,unicode-math}
\setmathfont{Latin Modern Math}
\usepackage{enumitem,needspace,xurl,hyperref}
\setlength{\parindent}{0pt}
\setlength{\parskip}{5pt}
\setlength{\emergencystretch}{3em}
\newcommand{\N}{\mathbb{N}}
\newcommand{\Z}{\mathbb{Z}}
\newcommand{\Q}{\mathbb{Q}}
\newcommand{\R}{\mathbb{R}}
\newcommand{\C}{\mathbb{C}}
\newcommand{\F}{\mathbb{F}}
\newcommand{\A}{\mathbb{A}}
\newcommand{\PP}{\mathbb{P}}
\newcommand{\E}{\mathbb{E}}
\newcommand{\eps}{\varepsilon}
\newcommand{\dd}{\,\mathrm{d}}
\begin{document}
\section*{UnsolvedMath ID: 2791}
\textbf{Source problem code:} catalogue entry
\subsection{Definitions and mathematical statement}
Let $X\subset\R^2$ be compact and totally disconnected, with $|X|\geq2$. The source denotes by $\Gamma_X$ the mapping class group of $\R^2\setminus X$. Its plane actions, end restrictions, and orientation-preserving circle actions suggest the convention of orientation-preserving plane homeomorphisms preserving $X$, modulo isotopies of that type, with the end at infinity distinguished. Those orientation and end conventions are not explicitly stated and must be fixed; the unrestricted mapping class group of the punctured surface need not have the same end action. Let $P\Gamma_X$ denote the pure subgroup acting trivially on $X$.

The source asks five distinct questions.
\begin{enumerate}
\item For $Y\subset X$, let $\Gamma_{X,Y}$ be the subgroup preserving $Y$ setwise, with restriction or forgetful homomorphism $r_{X,Y}:\Gamma_{X,Y}\to\Gamma_Y$ when this is defined. Classify families of real-valued maps $q_X:\Gamma_X\to\R$ satisfying
\[
\forall g\in\Gamma_{X,Y},\qquad q_X(g)\geq q_Y(r_{X,Y}(g)).
\]
This inequality, not equality, is the source definition of functoriality. No additivity, homogeneity, or complete conjugacy-detection property is imposed by that definition.
\item For a plane homeomorphism $\alpha$ preserving $X$, let $C_\alpha(X)$ consist of $\alpha$-invariant closed subsets $Y$ of $X$, with Hausdorff topology. Classify functorial $q$ for which $Y\mapsto q_Y([\alpha]_Y)$ is continuous or semicontinuous. On nonempty compact subsets, the Hausdorff metric is
\[
d_H(Y,Z)=\max\{\sup_{y\in Y}\inf_{z\in Z}|y-z|,\ \sup_{z\in Z}\inf_{y\in Y}|z-y|\}.
\]
Lower semicontinuity means $F(Y)\leq\liminf_jF(Y_j)$, and upper semicontinuity means $F(Y)\geq\limsup_jF(Y_j)$ for Hausdorff-convergent sequences; the source does not select a direction. Empty subsets and subsets with fewer than two points require additional domain conventions for $q_Y$.
\item Relate ray-graph translation length $\tau_X(g)$ to topological entropy of a representative $\alpha$. The displayed target is a positive lower comparison $h_{\mathrm{top}}(\alpha)\geq C\tau_X(g)$, where $g=[\alpha]_X$. A conventional stable length is
\[
\tau_X(g)=\lim_{m\to\infty}\frac{d(v,g^mv)}m
\]
for an isometric graph action, independent of the base vertex. The source does not say whether stable length or minimum displacement is intended, whether $C$ may depend on the fixed $X$, or which entropy definition on the noncompact plane is used. Entropy of the extension to the one-point compactification is one possible convention, not silently chosen here.
\item When $X$ is a planar Cantor set, let $N\leq P\Gamma_X$ be normal in $\Gamma_X$. For each positive integer $n$, restrict $N$ to an $n$-element subset $F\subset X$ and identify the resulting image with a subgroup $N_n$ of the pure $n$-puncture group $P\Gamma_n$. Determine which whole sequences $(N_n)_{n\geq1}$ occur. Identifications of finite subsets require a labeling or conjugacy convention; the word ``any'' in the source must be understood with that equivalence made explicit. The question concerns joint compatibility, not just the individual possible images.
\item For which $g\in\Gamma_X$ does some embedding of $X$ in the plane allow a representative by a $C^\infty$ diffeomorphism, or by a $C^k$ diffeomorphism for fixed finite $k$? The transport of $g$ under a changed embedding and the differentiability/ambient equivalence conventions must be specified. The source further asks whether realization or nonrealization can be detected by numerical data from ray-graph or circle actions.
\end{enumerate}

The supplied ray graph $\mathcal R_X$ has vertices isotopy classes of rays from infinity to $X$ and edges given by disjoint realizations. The notes report connectedness, hyperbolicity, infinite diameter, contractibility of its flag complex, and an isometric $\Gamma_X$ action in the intended setting. They give a cyclic completion circle $S_X^1$ with orientation-preserving action, and examples of numerical invariants: rotation number in $\R/\Z$, translation length, and path-orbit-dependent counting quasimorphisms. Rotation number is not itself a real-valued $q$ without an additional convention. The notes report an inclusion-induced map $\mathcal R_X\to\mathcal R_Y$ that is $1$-Lipschitz up to bounded-choice ambiguity, explaining the translation-length monotonicity. Full hypotheses for these ray-graph assertions, especially the finite small-$X$ cases permitted initially, are not supplied and are not assumed uniformly for every such $X$.

For the source braid-forcing supplement, $g<h$ means every disk homeomorphism representing $g$ on some finite invariant set also represents $h$ on another invariant set. With $G_g$ the set of representatives, the supplied definition is $G_g\subseteq G_h$. It asks about the relation of this order to continuity and semicontinuity, and extends the question to all $\Gamma_X$. The notes propose suprema of functions over $G_g$ as monotone quantities; an arbitrary real-valued function can have infinite supremum, so boundedness or extended-value conventions are needed for that construction to stay real-valued. The notes also report that Calegari--Chen identify individual groups $N_n$ by an ``inertia'' condition, whose definition and sequence-compatibility data are not supplied.

The registry is open and the literature triage dated 2026-08-17 reports partial ray-graph progress but no general classification for this multi-part programme. The mathematical questions and every supplementary scope are retained rather than reduced to only the first two invariant questions. All reported structures, results, and statuses are supplied information, not independently checked theorems here.
\subsection{Short English statement}
Classify functorial and continuity-compatible invariants for planar totally disconnected-end mapping class groups, relate ray-graph dynamics to entropy, classify finite restriction images of normal pure subgroups, and characterize smooth realizability after embedding changes. The source also asks about braid forcing; several domain, topology, and numerical conventions remain unspecified.
\subsection{Sources}
\begin{itemize}
\item K3: A New Problem List in Low-Dimensional Topology, Problem 2.43 (2026), supplied Kirby-style source: \url{https://aimath.org/pastworkshops/kirbylistrep.pdf}. This citation is retained without opening or generating a PDF.
\item [Boy88] Philip Boyland, ``An analog of Sharkovski's theorem for twist maps,'' Hamiltonian dynamical systems (Boulder, CO, 1987), Contemporary Mathematics 81 (1988), 119--133, American Mathematical Society, Providence, RI, DOI 10.1090/conm/081/986261.
\item [CC22] Danny Calegari and Lvzhou Chen, ``Normal subgroups of big mapping class groups,'' Transactions of the American Mathematical Society, Series B 9 (2022), 957--976, DOI 10.1090/btran/108.
\item Local registry data/unsolvedmath/problems.json, permanent catalogue entry, source code catalogue entry, ``Kirby Problem 2.43''; all five questions, complete supplementary notes, and literature triage dated 2026-08-17. Dataset attribution: Ulam AI and the UnsolvedMath Contributors.
\end{itemize}
\end{document}
